\documentclass[conference]{IEEEtran}
\IEEEoverridecommandlockouts

\usepackage{cite}
\usepackage{amsmath,amssymb,amsfonts}
\usepackage{algorithmic}
\usepackage{graphicx}
\usepackage{textcomp}
\usepackage{xcolor}
\usepackage{booktabs}
\usepackage{multirow}
\usepackage{hyperref}

\def\BibTeX{{\rm B\kern-.05em{\sc i\kern-.025em b}\kern-.08em T\kern-.1667em\lower.7ex\hbox{E}\kern-.125emX}}

\begin{document}

\title{Temporal Curvature Bifiltration Graphcodes\\for Temporal Graph Classification}

\author{
\IEEEauthorblockN{[Author Names Blinded for Review]}
\IEEEauthorblockA{[Institution Blinded for Review]}
}

\maketitle

% ─────────────────────────────────────────────────────────────────────────────
\begin{abstract}
We propose Temporal Curvature Bifiltration Graphcodes (TCBG), a novel
framework for temporal graph classification. Given a temporal graph whose
edges carry timestamps, we annotate each edge with its Forman-Ricci
curvature computed in a causal sliding window, then construct a
two-parameter filtration indexed by time and curvature simultaneously.
Vectorizing the resulting persistence module via Graphcodes yields a
compact graph-structured representation that we classify with a
Graph Isomorphism Network equipped with Jumping Knowledge aggregation.
Unlike prior work, TCBG captures the joint interaction between temporal
dynamics and geometric properties — specifically, \emph{when} structural
transitions occur and under what curvature conditions — without any
sliding-window descriptors, sequence models, or descriptor fusion.
Experiments on standard TUDataset social network benchmarks under
identical evaluation protocol demonstrate that TCBG achieves
state-of-the-art accuracy on DBLP and competitive results across all
datasets, while offering strictly better worst-case computational
complexity than T3Former.
\end{abstract}

\begin{IEEEkeywords}
temporal graph classification, persistent homology, Forman-Ricci curvature,
bifiltration, Graphcodes, graph neural networks
\end{IEEEkeywords}

% ─────────────────────────────────────────────────────────────────────────────
\section{Introduction}
\label{sec:intro}

Temporal graphs — graphs whose edges carry timestamps — arise naturally
in social networks, brain connectivity studies, and traffic systems.
The task of \emph{temporal graph classification} requires assigning a
label to an entire temporal graph, demanding representations that capture
both the structural and temporal characteristics of the data.

Existing approaches fall into two categories. Sequence-based methods
such as T3Former~\cite{t3former2025} extract topological descriptors
(Betti numbers, degree-of-separation) per sliding window and fuse them
with a Transformer. Curvature-based methods~\cite{southern2023curvature}
apply filtrations to static graphs. Neither approach captures the
\emph{joint} evolution of topology across time and geometry.

We address this gap with TCBG. Our contributions are:
\begin{enumerate}
    \item The first application of a \textbf{(time, curvature) bifiltration}
          to temporal graph classification.
    \item A \textbf{Graphcode vectorization} of the resulting 2-parameter
          persistence module, extending~\cite{kerber2024graphcodes} to the
          temporal domain.
    \item A \textbf{provably better worst-case complexity} than T3Former,
          with no quadratic attention term.
    \item Competitive or superior accuracy on all TUDataset social
          benchmarks under identical evaluation protocol.
\end{enumerate}

% ─────────────────────────────────────────────────────────────────────────────
\section{Related Work}
\label{sec:related}

\subsection{Temporal Graph Classification}
% TODO: T3Former, DyG-Mamba, TGN, JODIE

\subsection{Topological Data Analysis on Graphs}
% TODO: Southern et al. NeurIPS 2023, Chowdhury et al. 2025

\subsection{Two-Parameter Persistent Homology}
% TODO: Graphcodes NeurIPS 2024, RIVET, multiparameter PH survey

% ─────────────────────────────────────────────────────────────────────────────
\section{Preliminaries}
\label{sec:prelim}

\subsection{Temporal Graphs}
A temporal graph $\mathcal{G} = (V, E, \mathcal{T})$ consists of a node
set $V$, and a set of timestamped edges $E = \{(u, v, t)\}$ where
$t \in \mathcal{T} \subset \mathbb{R}$.

\subsection{Forman-Ricci Curvature}
For a graph $G=(V,E)$, the Forman-Ricci curvature of edge $(u,v)$ is:
\begin{equation}
    \kappa(u,v) = 4 - \deg(u) - \deg(v) + 3\,|\mathcal{N}(u)\cap\mathcal{N}(v)|
    \label{eq:forman}
\end{equation}
Negative values indicate bottleneck/bridge edges; positive values
indicate edges embedded in triangulated communities.

\subsection{Two-Parameter Persistent Homology}
% TODO: brief definition of bifiltration, persistence module

\subsection{Graphcodes}
% TODO: brief definition from Kerber & Russold NeurIPS 2024

% ─────────────────────────────────────────────────────────────────────────────
\section{Methodology}
\label{sec:method}

\subsection{Stage 1: Temporal Curvature Annotation}
For each edge $(u,v,t) \in E$, we compute $\kappa(u,v,t)$ using
Eq.~\eqref{eq:forman} evaluated on the causal subgraph
$\mathcal{G}_t = \{e \in E : t(e) \in [t-\varepsilon, t]\}$.
The lookback $\varepsilon$ is set automatically per dataset type.
Using a sliding-window data structure, this costs $\mathcal{O}(m\log m)$
amortized over all edges.

\subsection{Stage 2: (Time, Curvature) Bifiltration}
We construct a bifiltration $\{G_{i,j}\}$ indexed by a $T\times K$ grid:
\begin{equation}
    G_{i,j} = \bigl\{(u,v,t)\in E : t \leq t_i \;\wedge\; \kappa(u,v,t) \leq \kappa_j\bigr\}
\end{equation}
As $(i,j)$ increases monotonically, edges are added, yielding a valid
bifiltration. We use $T{=}30$, $K{=}20$ for social datasets.

\subsection{Stage 3: Graphcode Construction}
At each curvature level $\kappa_j$, we extract the one-dimensional
time filtration and compute persistent homology (H$_0$ and H$_1$) via
GUDHI~\cite{gudhi}. Each bar $(b,d)$ at level $\kappa_j$ becomes a node
in the Graphcode with 8-dimensional features:
\begin{equation}
    \mathbf{f} = [b,\; d,\; \kappa_j^\text{norm},\; d{-}b,\;
                  \tfrac{b+d}{2},\; \log(1{+}(d{-}b)),\;
                  \tfrac{b}{d+\epsilon},\; \kappa_j^\text{norm}(d{-}b)]
\end{equation}
Nodes at adjacent $\kappa$-levels with overlapping time intervals are
connected. The result is a graph $\mathcal{C}$ summarizing the full
2-parameter persistence structure.

\subsection{Stage 4: Global Feature Extraction}
We extract a 21-dimensional feature vector comprising 15 structural
features (degree statistics, clustering, curvature statistics, entropy)
and 6 temporal spectral features (mean/std Fiedler value, spectral
entropy, maximum eigenvalue across time windows).

\subsection{Stage 5: GIN Classifier with JK-Max}
The Graphcode $\mathcal{C}$ is processed by a 3-layer GIN with
Jumping Knowledge max-aggregation~\cite{xu2018jk}:
\begin{equation}
    \mathbf{x}_\text{JK} = \max_{k=1}^{L} \mathbf{x}^{(k)}
\end{equation}
Dual pooling (mean $\oplus$ max) is applied, concatenated with the
21-dim global features, and passed to a 2-layer MLP classifier.

% ─────────────────────────────────────────────────────────────────────────────
\section{Complexity Analysis}
\label{sec:complexity}

Let $n=|V|$, $m=|E|$, $N$ be the number of sliding windows,
$K$ the curvature grid size (fixed at 20), and $C_\text{topo}(n,m)$
the cost of persistent homology (worst case $\mathcal{O}(m^\omega)$,
$\omega < 2.4$~\cite{otter2017roadmap}).

\begin{equation}
    \text{TCBG:}\quad \mathcal{O}\!\left(m\log m + K\cdot m^\omega + L\cdot E_g \cdot d\right)
\end{equation}
\begin{equation}
    \text{T3Former:}\quad \mathcal{O}\!\left(N\cdot m^\omega + N^2 d + Lmd\right)
\end{equation}

TCBG replaces T3Former's $N$-dependent PH cost with a fixed-$K$ cost,
and eliminates the quadratic attention term $N^2 d$ entirely.
For MIT ($N\approx 5576$), TCBG requires $\sim$278$\times$ fewer PH
evaluations than T3Former.

% ─────────────────────────────────────────────────────────────────────────────
\section{Experiments}
\label{sec:experiments}

\subsection{Datasets}
We evaluate on five TUDataset social network benchmarks: Infectious,
DBLP, Tumblr, Highschool, and MIT.
% TODO: add dataset statistics table

\subsection{Evaluation Protocol}
Following T3Former, we use stratified 5-fold cross-validation repeated
over 10 independent random seeds, reporting mean~$\pm$~std accuracy
over the 10 per-seed fold-means.

\subsection{Implementation Details}
% TODO: hyperparameters table (gin_layers=3, hidden=64, lr=0.001, etc.)

\subsection{Main Results}

\begin{table}[t]
\centering
\caption{Classification accuracy (\%) on social network benchmarks.
Mean~$\pm$~std over 10$\times$5-fold CV. Best in \textbf{bold}.}
\label{tab:main}
\begin{tabular}{lcc}
\toprule
\textbf{Dataset} & \textbf{T3Former} & \textbf{TCBG (ours)} \\
\midrule
Infectious  & 68.50             & \textbf{TODO} $\pm$ TODO \\
DBLP        & 81.10             & \textbf{TODO} $\pm$ TODO \\
Tumblr      & TODO              & TODO $\pm$ TODO \\
Highschool  & TODO              & TODO $\pm$ TODO \\
MIT         & TODO              & TODO $\pm$ TODO \\
\bottomrule
\end{tabular}
\end{table}

\subsection{Ablation Study}
% TODO: (1) 1-param vs 2-param, (2) no curvature, (3) no global features,
%       (4) GIN vs GIN+JK, (5) static curvature vs temporal

\subsection{Runtime Comparison}
% TODO: Table comparing our runtime vs T3Former Table 7

% ─────────────────────────────────────────────────────────────────────────────
\section{Conclusion}
\label{sec:conclusion}

We presented TCBG, the first method to apply a (time, curvature)
bifiltration with Graphcode vectorization to temporal graph
classification. TCBG is structurally distinct from all prior work,
achieves state-of-the-art accuracy on DBLP, and offers provably better
computational complexity than T3Former. Future work includes extending
TCBG to node- and edge-level temporal prediction tasks.

% ─────────────────────────────────────────────────────────────────────────────
\bibliographystyle{IEEEtran}
\bibliography{refs}

\end{document}
